\section{How we built the corpus}
\label{sec:corpus}

\subsection{Discovery}

We started from 51 seed papers found during scoping, 40 of them directly
relevant to the two tracks and 11 related surveys kept for comparison. We
queried the Semantic Scholar Academic Graph API and the arXiv API with 14
structured keyword clusters, seven per track, drawn from the established terms
of each literature. That added 5,069 unique candidates after deduplication on
DOI, then arXiv identifier, then fuzzy title and author match.

From the seeds we ran one hop of backward and forward snowballing through the
Semantic Scholar citation graph. This is the step that matters for a
two-literature study: keyword search finds what a community calls its own work,
and snowballing finds what it reads. Snowballing added 20,801 further
candidates. The discovery pool reached 25,922 papers with 45,861 recorded
citation relationships.

Querying both tracks with one protocol is the methodological move this study
needs. No prior systematization in this space runs a single search across both
communities, so none can measure the distance between them.

\subsection{Screening}

Dual coding of 25,922 candidates is not feasible for a small team, so we
applied a precision-oriented filter. A candidate was kept for manual review if
it was one of the 51 seeds, or if it passed an automated inclusion screen and
named a recognized method or author from either literature. The inclusion
screen required language indicating an LLM agent with tool use, retrieval,
memory or autonomous operation. It excluded pure jailbreak work and
training-time-only poisoning.

This yielded 1,158 papers for manual screening, of which 984 were included and
174 excluded. Every search query, timestamp, hit count and raw API response was
logged, and every automated decision is retained next to any later human
correction, so the screening history stays auditable.

\subsection{Full-text coding}

All 984 included papers were read in full, from abstract to conclusion, not by
title and abstract alone. Each paper carries 32 recorded columns, of which
\textbf{17 are coded from the full text} and the rest are bibliographic
metadata retrieved from the search APIs. The coded fields are track, screening
decision, channel, consequence, defense intervention point, temporal
persistence, the two cross-citation fields, evidence grade, artifact release,
models evaluated, datasets and benchmarks, key result, baselines compared,
threat model, stated limitations, and a technical summary.

Two coding choices are worth stating because they change the numbers.

First, \textbf{the cross-citation fields are extracted from the full text, not
derived from the citation graph}. A paper counts as citing the other track when
its reference list contains a paper that resolves to that track in our graph.
We resolve this mechanically against every cited paper in the full discovery
pool, not just against the included corpus, so the measurement is not an
artifact of what we chose to include.

Second, \textbf{each defense is coded by what its own paper actually evaluated},
not by what it claims to address. A defense that says it addresses context
corruption in general but evaluates only on knowledge-conflict benchmarks is
coded as validated against the incidental threat model. This single choice
produces the central number in Section~\ref{sec:taxonomy}.

\subsection{Verification}

LLM-assisted extraction at this scale carries an error rate, and every count in
this paper depends on it. We therefore verified the coding with human readers
against the original PDFs. Verification is stratified: the fields that carry
the headline claims (track, channel, consequence, role, and the two
cross-citation fields) were verified in full, and the remaining fields were
verified on a random sample large enough to estimate an error rate.
Disagreements were recorded rather than silently overwritten, so inter-rater
agreement is reported rather than assumed.

\subsection{Citation network}

The citation network used in Section~\ref{sec:literatures} is rebuilt from the
PDFs rather than from API metadata. We extract each paper's reference section
and match references against corpus titles by a space-free contiguous prefix
that is unique within the corpus.

Rebuilding from PDFs rather than trusting API edges is not a detail. API
citation data is incomplete for recent preprints, which is most of this field.
Two implementation points cost us real edges before we found them. Matching
must ignore whitespace, because PDF extraction breaks hyphenated titles across
lines. Reference-section detection must accept headings other than the English
word "References", because venues in the corpus use
"Bibliographical References", small-capitals variants, and non-English
headings. Each of these, left unfixed, silently deletes true edges and would
have made the two literatures look even more separate than they are.
